\section{Программная реализация задачи}

\subsection{Структура программы}

Программа реализована на языке Java 17 с графическим интерфейсом на Swing. Структура проекта:

\begin{itemize}
\item \texttt{engine/} — численные методы интерполяции (чистые вычисления)
\item \texttt{engine/common/} — общие утилиты для центрально-разностных методов
\item \texttt{util/} — выбор подходящей формулы (вперёд/назад)
\item \texttt{model/} — модели данных (InterpolationResult, InputData)
\item \texttt{service/} — загрузка данных (из таблицы, файла, функции)
\item \texttt{gui/} — графический интерфейс (Swing)
\end{itemize}

\subsection{Конечные разности (FiniteDifferenceTable.java)}

\begin{lstlisting}[style=javaStyle, label=code:fd]
public class FiniteDifferenceTable {

    public static double[][] build(double[] y) {
        int n = y.length;
        double[][] table = new double[n][n];
        for (int i = 0; i < n; i++) {
            table[i][0] = y[i];
        }
        for (int j = 1; j < n; j++) {
            for (int i = 0; i < n - j; i++) {
                table[i][j] = table[i + 1][j - 1] - table[i][j - 1];
            }
        }
        return table;
    }
}
\end{lstlisting}

\subsection{Полином Лагранжа (LagrangeSolver.java)}

\begin{lstlisting}[style=javaStyle, label=code:lagrange]
public class LagrangeSolver {

    public static InterpolationResult interpolate(double[] x, double[] y,
                                                  double xTarget) {
        int n = x.length;
        double result = 0.0;

        for (int i = 0; i < n; i++) {
            double term = y[i];
            for (int j = 0; j < n; j++) {
                if (i != j) {
                    term *= (xTarget - x[j]) / (x[i] - x[j]);
                }
            }
            result += term;
        }

        double[][] diffTable = FiniteDifferenceTable.build(y);
        return new InterpolationResult(
                "\u041f\u043e\u043b\u0438\u043d\u043e\u043c \u041b\u0430\u0433\u0440\u0430\u043d\u0436\u0430",
                xTarget, result, diffTable, n - 1);
    }
}
\end{lstlisting}

\subsection{Полином Ньютона с разделёнными разностями (NewtonSolver.java)}

\begin{lstlisting}[style=javaStyle, label=code:newton]
public class NewtonSolver {

    private static double[][] buildDividedDiffTable(double[] x, double[] y) {
        int n = y.length;
        double[][] table = new double[n][n];
        for (int i = 0; i < n; i++) {
            table[i][0] = y[i];
        }
        for (int j = 1; j < n; j++) {
            for (int i = 0; i < n - j; i++) {
                table[i][j] = (table[i + 1][j - 1] - table[i][j - 1])
                        / (x[i + j] - x[i]);
            }
        }
        return table;
    }

    public static InterpolationResult interpolateForward(double[] x, double[] y,
                                                         double xTarget) {
        double[][] dd = buildDividedDiffTable(x, y);
        int n = x.length;
        double result = dd[0][0];
        double product = 1.0;

        for (int i = 1; i < n; i++) {
            product *= (xTarget - x[i - 1]);
            result += dd[0][i] * product;
        }

        double[][] fd = FiniteDifferenceTable.build(y);
        return new InterpolationResult(
                "Ньютон (вперёд)", xTarget, result, fd, n - 1);
    }

    public static InterpolationResult interpolateBackward(double[] x, double[] y,
                                                          double xTarget) {
        double[][] dd = buildDividedDiffTable(x, y);
        int n = x.length;
        double result = dd[n - 1][0];
        double product = 1.0;

        for (int i = 1; i < n; i++) {
            product *= (xTarget - x[n - i]);
            result += dd[n - 1 - i][i] * product;
        }

        double[][] fd = FiniteDifferenceTable.build(y);
        return new InterpolationResult(
                "Ньютон (назад)", xTarget, result, fd, n - 1);
    }
}
\end{lstlisting}

\subsection{Общие утилиты центрально-разностных методов (CentralDifferenceUtils.java)}

\begin{lstlisting}[style=javaStyle, label=code:cdutil]
public class CentralDifferenceUtils {

    public static double factorial(int k) {
        double f = 1.0;
        for (int i = 2; i <= k; i++) f *= i;
        return f;
    }

    public static double gaussForwardCoeff(double t, int k) {
        double prod = 1.0;
        if (k % 2 == 1) {
            int m = (k - 1) / 2;
            for (int s = -m; s <= m; s++) {
                prod *= (t + s);
            }
        } else {
            int start = -k / 2;
            int end = k / 2 - 1;
            for (int s = start; s <= end; s++) {
                prod *= (t + s);
            }
        }
        return prod / factorial(k);
    }

    public static double gaussBackwardCoeff(double t, int k) {
        double prod = 1.0;
        if (k % 2 == 1) {
            int m = (k - 1) / 2;
            for (int s = -m; s <= m; s++) {
                prod *= (t + s);
            }
        } else {
            int start = -(k / 2 - 1);
            int end = k / 2;
            for (int s = start; s <= end; s++) {
                prod *= (t + s);
            }
        }
        return prod / factorial(k);
    }
}
\end{lstlisting}

\subsection{Полином Гаусса (GaussSolver.java)}

\begin{lstlisting}[style=javaStyle, label=code:gauss]
public class GaussSolver {

    public static InterpolationResult interpolateForward(double[] x, double[] y,
                                                         double xTarget) {
        int n = y.length;
        double[][] ft = FiniteDifferenceTable.build(y);
        double h = x[1] - x[0];
        int c = n / 2;
        double t = (xTarget - x[c]) / h;
        double result = ft[c][0];

        for (int k = 1; k < n; k++) {
            double coeff = CentralDifferenceUtils.gaussForwardCoeff(t, k);
            int row = (k % 2 == 1) ? c - (k - 1) / 2 : c - k / 2;
            if (row >= 0 && row < n - k) {
                result += coeff * ft[row][k];
            }
        }

        return new InterpolationResult(
                "Гаусс 1-я формула (вперёд)", xTarget, result, ft, n - 1);
    }

    public static InterpolationResult interpolateBackward(double[] x, double[] y,
                                                          double xTarget) {
        int n = y.length;
        double[][] ft = FiniteDifferenceTable.build(y);
        double h = x[1] - x[0];
        int c = n / 2;
        double t = (xTarget - x[c]) / h;
        double result = ft[c][0];

        for (int k = 1; k < n; k++) {
            double coeff = CentralDifferenceUtils.gaussBackwardCoeff(t, k);
            int row = (k % 2 == 1) ? c - (k + 1) / 2 : c - k / 2;
            if (row >= 0 && row < n - k) {
                result += coeff * ft[row][k];
            }
        }

        return new InterpolationResult(
                "Гаусс 2-я формула (назад)", xTarget, result, ft, n - 1);
    }
}
\end{lstlisting}

\subsection{Формула Стирлинга (StirlingSolver.java)}

Формула Стирлинга реализована как среднее арифметическое первой и второй формул Гаусса:

\begin{lstlisting}[style=javaStyle, label=code:stirling]
public class StirlingSolver {

    public static InterpolationResult interpolate(double[] x, double[] y,
                                                  double xTarget) {
        int n = y.length;
        double[][] ft = FiniteDifferenceTable.build(y);
        double h = x[1] - x[0];
        int c = n / 2;
        double t = (xTarget - x[c]) / h;

        double gf = gaussForwardEval(ft, c, t, n);
        double gb = gaussBackwardEval(ft, c, t, n);
        double result = (gf + gb) / 2.0;

        return new InterpolationResult(
                "Формула Стирлинга", xTarget, result, ft, n - 1);
    }

    private static double gaussForwardEval(double[][] ft, int c,
                                           double t, int n) {
        double result = ft[c][0];
        for (int k = 1; k < n; k++) {
            double coeff = CentralDifferenceUtils.gaussForwardCoeff(t, k);
            int row = (k % 2 == 1) ? c - (k - 1) / 2 : c - k / 2;
            if (row >= 0 && row < n - k) {
                result += coeff * ft[row][k];
            }
        }
        return result;
    }

    private static double gaussBackwardEval(double[][] ft, int c,
                                            double t, int n) {
        double result = ft[c][0];
        for (int k = 1; k < n; k++) {
            double coeff = CentralDifferenceUtils.gaussBackwardCoeff(t, k);
            int row = (k % 2 == 1) ? c - (k + 1) / 2 : c - k / 2;
            if (row >= 0 && row < n - k) {
                result += coeff * ft[row][k];
            }
        }
        return result;
    }
}
\end{lstlisting}

\subsection{Формула Бесселя (BesselSolver.java)}

\begin{lstlisting}[style=javaStyle, label=code:bessel]
public class BesselSolver {

    public static InterpolationResult interpolate(double[] x, double[] y,
                                                  double xTarget) {
        int n = y.length;
        double[][] ft = FiniteDifferenceTable.build(y);
        double h = x[1] - x[0];
        int c = n / 2;
        double t = (xTarget - x[c]) / h;
        double result = ft[c][0];

        for (int k = 1; k < n; k++) {
            double coeff = besselCoeff(t, k);
            if (Math.abs(coeff) < 1e-15) continue;

            double diffVal;
            if (k % 2 == 0) {
                int row1 = c - k / 2;
                int row2 = c - k / 2 + 1;
                if (row1 >= 0 && row1 < n - k && row2 >= 0 && row2 < n - k) {
                    diffVal = (ft[row1][k] + ft[row2][k]) / 2.0;
                } else if (row1 >= 0 && row1 < n - k) {
                    diffVal = ft[row1][k];
                } else if (row2 >= 0 && row2 < n - k) {
                    diffVal = ft[row2][k];
                } else {
                    continue;
                }
            } else {
                int row = c - (k - 1) / 2;
                if (row >= 0 && row < n - k) {
                    diffVal = ft[row][k];
                } else {
                    continue;
                }
            }
            result += coeff * diffVal;
        }

        return new InterpolationResult(
                "Формула Бесселя", xTarget, result, ft, n - 1);
    }

    private static double besselCoeff(double t, int k) {
        if (k == 1) return t;

        double prod = 1.0;
        if (k % 2 == 0) {
            int m = k / 2;
            for (int i = -(m - 1); i <= m - 1; i++) {
                prod *= (t + i);
            }
            prod *= (t - m);
        } else {
            int m = (k - 1) / 2;
            prod *= (t - 0.5);
            for (int i = -(m - 1); i <= m - 1; i++) {
                prod *= (t + i);
            }
            prod *= (t - m);
        }
        return prod / CentralDifferenceUtils.factorial(k);
    }
}
\end{lstlisting}

\subsection{Графический интерфейс (MainFrame.java)}

Приложение реализует три способа ввода данных:
\begin{enumerate}
\item \textbf{Ручной ввод} — ввод значений $x$ и $y$ через запятую
\item \textbf{Загрузка из файла} — чтение табличных данных из текстового файла
\item \textbf{Генерация по функции} — выбор функции ($\sin x$, $\cos x$, $e^x$, $x^2$, $\ln x$), интервала и количества точек
\end{enumerate}

Пользователь может выбрать произвольный набор методов интерполяции (Лагранж, Ньютон, Гаусс, Стирлинг, Бессель) и указать два значения аргумента $X_1$ и $X_2$.

Результаты отображаются в виде:
\begin{itemize}
\item Таблицы конечных разностей
\item Вычисленных значений для каждого метода
\item Графика с узлами интерполяции, интерполяционной кривой и отмеченными точками $X_1$, $X_2$
\end{itemize}

\section{Пример работы программы}

\begin{figure}[H]
    \centering
    \includegraphics[width=1\linewidth]{pic/screenshot_main.png}
    \caption{Основное окно программы с данными варианта 14}
    \label{fig:main}
\end{figure}

\begin{figure}[H]
    \centering
    \includegraphics[width=1\linewidth]{pic/screenshot_function.png}
    \caption{Интерполяция функции $\sin(x)$}
    \label{fig:func}
\end{figure}

\section{Анализ результатов}

\subsection{Проверка корректности}

Все методы интерполяции (Лагранж, Ньютон, Гаусс, Стирлинг, Бессель) дают одинаковые значения функции для одних и тех же узлов и аргументов. Это объясняется тем, что интерполяционный полином степени $n-1$ для $n$ узлов единственен — различные формулы представляют собой лишь разные формы записи одного и того же полинома.

\subsection{Сравнение результатов}

Для варианта 14 (таблица 1.4) получены следующие результаты:
\begin{itemize}
\item $f(1{,}112) \approx 0{,}749956$ (все методы)
\item $f(1{,}319) \approx 2{,}845036$ (все методы)
\end{itemize}

\subsection{Тестирование на некорректных данных}

Программа корректно обрабатывает:
\begin{itemize}
\item Несовпадение количества значений $x$ и $y$
\item Немонотонные значения $x$
\item Пустые данные
\item Некорректные числовые форматы
\end{itemize}

\subsection{Выводы}

\begin{enumerate}
\item Все реализованные методы интерполяции дают идентичные результаты при использовании полного набора конечных разностей
\item Для интерполяции вблизи начала таблицы целесообразно использовать первую формулу Ньютона
\item Для интерполяции вблизи середины таблицы удобны формулы Гаусса, Стирлинга и Бесселя
\item Формула Стирлинга даёт наилучшие результаты при $|q| \leq 0{,}25$
\item Формула Бесселя оптимальна при $|q| \approx 0{,}5$
\item Программная реализация с графическим интерфейсом позволяет наглядно сравнить результаты различных методов
\end{enumerate}

\section{Вывод}

В ходе выполнения лабораторной работы были изучены и реализованы основные методы интерполяции функций: полином Лагранжа, полином Ньютона с разделёнными разностями, полином Гаусса (первая и вторая формулы), а также дополнительные методы — формулы Стирлинга и Бесселя.

Вычислительная часть работы выполнена для варианта 14 (таблица 1.4): построена таблица конечных разностей, вычислены значения функции для $X_1 = 1{,}112$ и $X_2 = 1{,}319$ с использованием формул Ньютона, Гаусса и Лагранжа.

Программная часть реализована на языке Java с графическим интерфейсом Swing. Приложение поддерживает три способа ввода данных, позволяет выбрать набор методов интерполяции, отображает таблицу конечных разностей и строит график интерполяции. Программа протестирована на различных наборах данных, включая некорректные.

\end{enumerate}
